\section{What the therapist's turns do to the patient}
\label{sec:patient}

MI judges a therapist behavior by what the patient says next.
% \dnotedone{renamed: utterance coder, defined in §3.3}{[positional codes ?]}
% \dnotedone{answered in §3.3: each judge also runs the utterance coder}{not clear, i thought the judges were for something else and here you have a MITI classifier - /}
The codes are in transcript order, so each patient utterance can be followed through the
therapist's reply to the next patient utterance (Appendix Table~\ref{tab:process} at the
checkpoints of Table~\ref{tab:endpoint}; every iteration: Table~\ref{tab:process-all}).

% 2026-10-07 (step 16, Lior): the process table (tab:process) moved to Appendix A.


% \dnotedone{paragraph rewritten: one judge; K=0, K=5 and Base named}{very confusing. not clear when you are talking above 0 vs look ahead and base vs both..}
\textbf{How the therapist replies.}
% \dnotedone{endpoint dropped; the opening says iteration 10}{[endpoint ?]}
% \dnotedone{answered: reflection defined in §3.1, reply rates in Appendix D.3}{[reflects it?]}
% \dnotedone{metaphor dropped: now change talk, defined in §3.1}{[own case for change -?]}
When the patient expresses change talk, the $\Kf$ policy reflects it in $26\%$ of cases, against
$0.3\%$ for the $\Kz$ policy (Base $14\%$; Appendix Figure~\ref{fig:responsiveness}). A reflection
is also one of MI's answers to sustain talk \citep{miller2013mi}, and the $\Kf$ policy gives it
somewhat more often ($37\%$, against $26\%$ under $\Kz$ and $22\%$ at the Base), significantly
more than $\Kz$'s best checkpoint but not than its last. Each final policy also has its own MI-inconsistent answer
to sustain talk.
% \dnotedone{rephrased: praise regardless became praise in reply to sustain talk}{of what?}
% \dnotedone{rephrased: the reply of Table 2, measured across conversations}{[This is Table \ref{tab:excerpt} as a rate -?]}
% \dnotedone{renamed: K=0 / K=5 only in the results}{[turn-level policy - 0? use consistent names]}
The $\Kz$ policy praises: $32\%$ of its replies to sustain talk, against $2\%$ at the Base. The
$\Kf$ policy persuades: $39\%$ of its replies to sustain talk, against $24\%$ at the Base.
% \dnotedone{answered: praise premium measured, Appendix C.4}{[which the rubric rewards in the turn where it happens ??]}
% 2026-10-07 (step 16): the praise-reward numbers moved to the Discussion, which states them once.

% \dnotedone{cut; measured change-talk persistence replaces it}{[ pays in the turns the rubric is not shown at $K{=}0$. ?]}
\textbf{Whether change talk continues.} In the average conversation under the $\Kf$ policy, a
patient who has just expressed change talk expresses it again in their next utterance in $97\%$ of
cases, against $80\%$ under $\Kz$ ($88\%$ at its best checkpoint; Base $75\%$;
Figure~\ref{fig:process}d).
% \dnotedone{cut: reading by therapist code retired; reason given here}{[the other way round: ?]}
% \dnotedone{cut with the praise-yield paragraph}{[for the reason the next paragraph gives. \textasciitilde{}]}
No single kind of reply accounts for it: with the patient's previous utterance held fixed, the
$\Kf$ policy's complex reflections are followed by change talk about as often as its replies overall
(Appendix~\ref{app:tables}). After sustain talk, the next patient utterance is change talk in $33\%$ of cases
under $\Kf$, against $28\%$ under $\Kz$ and at the Base, a difference significant against $\Kz$'s
best checkpoint ($22\%$) but not against its last.

\textbf{The session as a whole.}
% \dnotedone{rephrased: iteration 10, one judge, each run named}{[the endpoint ($0.68$ vs.\ $0.53$ under the training oracle, $0.72$ vs.\ $0.55$ held out; $\dz\,0.57$ / $0.70$) - not clear]}
At iteration 10, $68\%$ of the patient's utterances are change talk under $\Kf$, against $53\%$
under $\Kz$ ($52\%$ at its best checkpoint; Base $45\%$), while the share of sessions that reach any
change talk at all does not differ ($92\%$ vs.\ $86\%$).
% \dnotedone{rephrased: the difference between the runs}{between waht}
% \dnotedone{cut: the praise-spiral clause dropped}{[at scale ?]}
% 2026-10-07 (review round 4): the late-session rise/fall (old Figure 3d) left with that panel.
The $\Kf$ policy's patients also write longer utterances and use a crude cue of disengagement
slightly more often (Appendix~\ref{app:tables}); with a simulated patient we cannot tell engagement
from verbosity.
% \dnotedone{rephrased: why neither is interpreted is now stated}{[rather than resolve. ?]}

The held-out judge agrees on the direction of every contrast at iteration 10
(Appendix Table~\ref{tab:process-heldout}), with the $\Kf$ policy persuading after $58\%$ of the
patient's sustain talk (Base $25\%$).
